% ============================================================
%  gamma_amplification.tex
%  Fragment: include in a parent document via % ============================================================
%  gamma_amplification.tex
%  Fragment: include in a parent document via % ============================================================
%  gamma_amplification.tex
%  Fragment: include in a parent document via % ============================================================
%  gamma_amplification.tex
%  Fragment: include in a parent document via \input{gamma_amplification}
%
%  Required packages in parent preamble:
%    \usepackage{amsmath, amssymb, booktabs, array, bm}
%
%  Describes the fundamental gamma^2 amplification systematic
%  affecting kaon momentum reconstruction from decay-time beta
%  measurement in the KLong simulation.
% ============================================================

\subsection{Systematic Momentum Underestimation: the $\gamma^2$ Amplification Problem}

\subsubsection{The Momentum Formula}

The kaon momentum is reconstructed by measuring its velocity $\beta = v/c$ from the
flight path $L$ and the kaon decay time $\Delta t$:
\begin{equation}
  \beta_{\text{reco}} = \frac{L}{c\,\Delta t},
  \label{eq:beta_reco}
\end{equation}
and then inverting the relativistic energy–momentum relation to obtain
\begin{equation}
  p = m_K\,\gamma\,\beta\,c = \frac{m_K\,\beta\,c}{\sqrt{1-\beta^2}},
  \label{eq:p_from_beta}
\end{equation}
where $m_K = 497.6\;\text{MeV}/c^2$ is the kaon mass.

\subsubsection{Error Propagation and the $\gamma^2$ Factor}

Differentiating Eq.~\eqref{eq:p_from_beta} with respect to $\beta$:
\begin{equation}
  \frac{\mathrm{d}p}{\mathrm{d}\beta}
  = m_K c\,\frac{\mathrm{d}}{\mathrm{d}\beta}\!\left[\frac{\beta}{\sqrt{1-\beta^2}}\right]
  = m_K c\,\frac{1}{(1-\beta^2)^{3/2}}
  = m_K c\,\gamma^3.
  \label{eq:dp_dbeta}
\end{equation}
Dividing by $p = m_K c\,\gamma\beta$ to form the relative error:
\begin{equation}
  \boxed{
  \frac{\delta p}{p} = \gamma^2\,\frac{\delta\beta}{\beta}.
  }
  \label{eq:gamma2_amplification}
\end{equation}
This is the central result.  The Lorentz factor squared $\gamma^2$ multiplies any
fractional error in $\beta$ to produce the fractional error in the reconstructed momentum.

\subsubsection{Physical Origin of the Amplification}

Near $\beta \to 1$ the function $p(\beta)$ becomes extremely steep
(Fig.~\ref{fig:p_vs_beta}).  A small absolute error $|\delta\beta|$ therefore maps to
an arbitrarily large absolute error in $p$.  Equivalently, a fixed relative tracking
uncertainty $\delta\beta/\beta$ is amplified by $\gamma^2$, which itself grows
without bound as $\beta \to 1$.

The amplification arises because the kaon is \emph{ultra-relativistic}: its total
energy $E \approx pc$, so a small fractional change in $\beta$ corresponds to a
small fractional change in $E$ but — through the subtraction
$p = \sqrt{E^2 - m_K^2 c^4}/c$ — to a much larger fractional change in $p$ whenever
$m_K c^2 \ll E$.

\subsubsection{Numerical Example — Event 891}

For the first clean reconstructed event in the
$T_{1}\text{-}400$ configuration simulation:

\begin{center}
\begin{tabular}{lrrl}
  \toprule
  Quantity & Value & Unit & Notes \\
  \midrule
  $p_{\text{true}}$  & 7.455  & GeV/$c$ & Monte~Carlo truth \\
  $\beta_{\text{true}}$ & 0.99773 & --- & from truth track \\
  $\gamma_{\text{true}}$ & 15.01  & --- & \\
  $\gamma^2_{\text{true}}$ & 225 & --- & amplification factor \\[4pt]
  $p_{\text{reco}}$  & 4.468  & GeV/$c$ & PoCA reconstruction \\
  $\delta p/p$       & $-40\%$ & --- & large underestimate \\[4pt]
  Implied $\delta\beta/\beta$ & $-0.18\%$ & --- &
        from Eq.~\eqref{eq:gamma2_amplification}: $-0.40/225$ \\
  \bottomrule
\end{tabular}
\end{center}

A sub-$0.2\%$ error in $\beta$ — well within detector resolution — produces a
$40\%$ underestimate of the momentum because $\gamma^2 = 225$.
Reducing tracking resolution cannot eliminate this effect; it can only marginally
reduce $\delta\beta$ while the $\gamma^2$ factor remains large.

\subsubsection{Identifying the Source of $\delta\beta$}

The measured $\beta$ depends on two independently reconstructed quantities:

\begin{equation}
  \beta_{\text{reco}} = \frac{L_{\text{reco}}}{c\,\Delta t_{\text{reco}}},
  \label{eq:beta_components}
\end{equation}

where
\begin{equation}
  L_{\text{reco}} = \bigl|\bm{x}_{\text{PoCA}} - \bm{x}_{\text{production}}\bigr|
  \label{eq:L_reco}
\end{equation}
is the flight-path length obtained from the PoCA vertex position $\bm{x}_{\text{PoCA}}$,
and $\Delta t_{\text{reco}}$ is derived from pion TOF measurements combined with
a back-propagation along the fitted pion tracks to the kaon production point.

Errors in either component contribute to $\delta\beta$:
\begin{equation}
  \frac{\delta\beta}{\beta} \approx \frac{\delta L}{L} - \frac{\delta(\Delta t)}{\Delta t}.
  \label{eq:beta_error_decomp}
\end{equation}

The PoCA algorithm minimises the separation between two straight-line tracks.
Even when the separation residual is small (e.g.\ $|\delta_{\text{PoCA}}| < 0.1\;\text{cm}$
for Event~891), the reconstructed $z$-coordinate of the vertex is sensitive to the
ratio of track slopes, and any bias in $z$ propagates directly into $L$ and hence
into $\beta$.

\subsubsection{Why Improving Track Resolution Helps Only Marginally}

Let $\sigma_\beta$ be the standard deviation of $\beta_{\text{reco}}$ arising from
finite track resolution.  The standard deviation of the reconstructed momentum is then:
\begin{equation}
  \sigma_p \approx m_K c\,\gamma^3\,\sigma_\beta
             = p\,\gamma^2\,\frac{\sigma_\beta}{\beta}.
  \label{eq:sigma_p}
\end{equation}
Halving $\sigma_\beta$ (e.g.\ by doubling the number of straw layers) halves
$\sigma_p$ — a modest gain — but the $\gamma^2$ prefactor is unchanged.  At
$p \gtrsim 5\;\text{GeV}/c$ where $\gamma^2 \gtrsim 100$, the achievable momentum
resolution is fundamentally limited by the decay-time method and the kaon mass, not
by the tracker hit resolution alone.

To improve sensitivity at high momentum one would need to either:
\begin{enumerate}
  \item measure the flight path $L$ independently (e.g.\ a displaced vertex tag), or
  \item measure the kaon energy directly (calorimetry).
\end{enumerate}
Neither is currently available in this detector geometry.

\subsubsection{Summary}

\begin{itemize}
  \item The relative momentum error scales as $\gamma^2 \times \delta\beta/\beta$
        (Eq.~\eqref{eq:gamma2_amplification}).
  \item At $p \sim 7.5\;\text{GeV}/c$, $\gamma^2 \approx 225$; a $0.18\%$ $\beta$
        error causes a $40\%$ momentum underestimate.
  \item Tracking resolution improvements reduce $\delta\beta$ but cannot remove
        the $\gamma^2$ amplification.
  \item The dominant error source — $z$-vertex bias or timing-chain bias — is
        isolated by the diagnostic reconstruction branches described in
        Section~\ref{sec:diagnostic_branches}.
\end{itemize}

%
%  Required packages in parent preamble:
%    \usepackage{amsmath, amssymb, booktabs, array, bm}
%
%  Describes the fundamental gamma^2 amplification systematic
%  affecting kaon momentum reconstruction from decay-time beta
%  measurement in the KLong simulation.
% ============================================================

\subsection{Systematic Momentum Underestimation: the $\gamma^2$ Amplification Problem}

\subsubsection{The Momentum Formula}

The kaon momentum is reconstructed by measuring its velocity $\beta = v/c$ from the
flight path $L$ and the kaon decay time $\Delta t$:
\begin{equation}
  \beta_{\text{reco}} = \frac{L}{c\,\Delta t},
  \label{eq:beta_reco}
\end{equation}
and then inverting the relativistic energy–momentum relation to obtain
\begin{equation}
  p = m_K\,\gamma\,\beta\,c = \frac{m_K\,\beta\,c}{\sqrt{1-\beta^2}},
  \label{eq:p_from_beta}
\end{equation}
where $m_K = 497.6\;\text{MeV}/c^2$ is the kaon mass.

\subsubsection{Error Propagation and the $\gamma^2$ Factor}

Differentiating Eq.~\eqref{eq:p_from_beta} with respect to $\beta$:
\begin{equation}
  \frac{\mathrm{d}p}{\mathrm{d}\beta}
  = m_K c\,\frac{\mathrm{d}}{\mathrm{d}\beta}\!\left[\frac{\beta}{\sqrt{1-\beta^2}}\right]
  = m_K c\,\frac{1}{(1-\beta^2)^{3/2}}
  = m_K c\,\gamma^3.
  \label{eq:dp_dbeta}
\end{equation}
Dividing by $p = m_K c\,\gamma\beta$ to form the relative error:
\begin{equation}
  \boxed{
  \frac{\delta p}{p} = \gamma^2\,\frac{\delta\beta}{\beta}.
  }
  \label{eq:gamma2_amplification}
\end{equation}
This is the central result.  The Lorentz factor squared $\gamma^2$ multiplies any
fractional error in $\beta$ to produce the fractional error in the reconstructed momentum.

\subsubsection{Physical Origin of the Amplification}

Near $\beta \to 1$ the function $p(\beta)$ becomes extremely steep
(Fig.~\ref{fig:p_vs_beta}).  A small absolute error $|\delta\beta|$ therefore maps to
an arbitrarily large absolute error in $p$.  Equivalently, a fixed relative tracking
uncertainty $\delta\beta/\beta$ is amplified by $\gamma^2$, which itself grows
without bound as $\beta \to 1$.

The amplification arises because the kaon is \emph{ultra-relativistic}: its total
energy $E \approx pc$, so a small fractional change in $\beta$ corresponds to a
small fractional change in $E$ but — through the subtraction
$p = \sqrt{E^2 - m_K^2 c^4}/c$ — to a much larger fractional change in $p$ whenever
$m_K c^2 \ll E$.

\subsubsection{Numerical Example — Event 891}

For the first clean reconstructed event in the
$T_{1}\text{-}400$ configuration simulation:

\begin{center}
\begin{tabular}{lrrl}
  \toprule
  Quantity & Value & Unit & Notes \\
  \midrule
  $p_{\text{true}}$  & 7.455  & GeV/$c$ & Monte~Carlo truth \\
  $\beta_{\text{true}}$ & 0.99773 & --- & from truth track \\
  $\gamma_{\text{true}}$ & 15.01  & --- & \\
  $\gamma^2_{\text{true}}$ & 225 & --- & amplification factor \\[4pt]
  $p_{\text{reco}}$  & 4.468  & GeV/$c$ & PoCA reconstruction \\
  $\delta p/p$       & $-40\%$ & --- & large underestimate \\[4pt]
  Implied $\delta\beta/\beta$ & $-0.18\%$ & --- &
        from Eq.~\eqref{eq:gamma2_amplification}: $-0.40/225$ \\
  \bottomrule
\end{tabular}
\end{center}

A sub-$0.2\%$ error in $\beta$ — well within detector resolution — produces a
$40\%$ underestimate of the momentum because $\gamma^2 = 225$.
Reducing tracking resolution cannot eliminate this effect; it can only marginally
reduce $\delta\beta$ while the $\gamma^2$ factor remains large.

\subsubsection{Identifying the Source of $\delta\beta$}

The measured $\beta$ depends on two independently reconstructed quantities:

\begin{equation}
  \beta_{\text{reco}} = \frac{L_{\text{reco}}}{c\,\Delta t_{\text{reco}}},
  \label{eq:beta_components}
\end{equation}

where
\begin{equation}
  L_{\text{reco}} = \bigl|\bm{x}_{\text{PoCA}} - \bm{x}_{\text{production}}\bigr|
  \label{eq:L_reco}
\end{equation}
is the flight-path length obtained from the PoCA vertex position $\bm{x}_{\text{PoCA}}$,
and $\Delta t_{\text{reco}}$ is derived from pion TOF measurements combined with
a back-propagation along the fitted pion tracks to the kaon production point.

Errors in either component contribute to $\delta\beta$:
\begin{equation}
  \frac{\delta\beta}{\beta} \approx \frac{\delta L}{L} - \frac{\delta(\Delta t)}{\Delta t}.
  \label{eq:beta_error_decomp}
\end{equation}

The PoCA algorithm minimises the separation between two straight-line tracks.
Even when the separation residual is small (e.g.\ $|\delta_{\text{PoCA}}| < 0.1\;\text{cm}$
for Event~891), the reconstructed $z$-coordinate of the vertex is sensitive to the
ratio of track slopes, and any bias in $z$ propagates directly into $L$ and hence
into $\beta$.

\subsubsection{Why Improving Track Resolution Helps Only Marginally}

Let $\sigma_\beta$ be the standard deviation of $\beta_{\text{reco}}$ arising from
finite track resolution.  The standard deviation of the reconstructed momentum is then:
\begin{equation}
  \sigma_p \approx m_K c\,\gamma^3\,\sigma_\beta
             = p\,\gamma^2\,\frac{\sigma_\beta}{\beta}.
  \label{eq:sigma_p}
\end{equation}
Halving $\sigma_\beta$ (e.g.\ by doubling the number of straw layers) halves
$\sigma_p$ — a modest gain — but the $\gamma^2$ prefactor is unchanged.  At
$p \gtrsim 5\;\text{GeV}/c$ where $\gamma^2 \gtrsim 100$, the achievable momentum
resolution is fundamentally limited by the decay-time method and the kaon mass, not
by the tracker hit resolution alone.

To improve sensitivity at high momentum one would need to either:
\begin{enumerate}
  \item measure the flight path $L$ independently (e.g.\ a displaced vertex tag), or
  \item measure the kaon energy directly (calorimetry).
\end{enumerate}
Neither is currently available in this detector geometry.

\subsubsection{Summary}

\begin{itemize}
  \item The relative momentum error scales as $\gamma^2 \times \delta\beta/\beta$
        (Eq.~\eqref{eq:gamma2_amplification}).
  \item At $p \sim 7.5\;\text{GeV}/c$, $\gamma^2 \approx 225$; a $0.18\%$ $\beta$
        error causes a $40\%$ momentum underestimate.
  \item Tracking resolution improvements reduce $\delta\beta$ but cannot remove
        the $\gamma^2$ amplification.
  \item The dominant error source — $z$-vertex bias or timing-chain bias — is
        isolated by the diagnostic reconstruction branches described in
        Section~\ref{sec:diagnostic_branches}.
\end{itemize}

%
%  Required packages in parent preamble:
%    \usepackage{amsmath, amssymb, booktabs, array, bm}
%
%  Describes the fundamental gamma^2 amplification systematic
%  affecting kaon momentum reconstruction from decay-time beta
%  measurement in the KLong simulation.
% ============================================================

\subsection{Systematic Momentum Underestimation: the $\gamma^2$ Amplification Problem}

\subsubsection{The Momentum Formula}

The kaon momentum is reconstructed by measuring its velocity $\beta = v/c$ from the
flight path $L$ and the kaon decay time $\Delta t$:
\begin{equation}
  \beta_{\text{reco}} = \frac{L}{c\,\Delta t},
  \label{eq:beta_reco}
\end{equation}
and then inverting the relativistic energy–momentum relation to obtain
\begin{equation}
  p = m_K\,\gamma\,\beta\,c = \frac{m_K\,\beta\,c}{\sqrt{1-\beta^2}},
  \label{eq:p_from_beta}
\end{equation}
where $m_K = 497.6\;\text{MeV}/c^2$ is the kaon mass.

\subsubsection{Error Propagation and the $\gamma^2$ Factor}

Differentiating Eq.~\eqref{eq:p_from_beta} with respect to $\beta$:
\begin{equation}
  \frac{\mathrm{d}p}{\mathrm{d}\beta}
  = m_K c\,\frac{\mathrm{d}}{\mathrm{d}\beta}\!\left[\frac{\beta}{\sqrt{1-\beta^2}}\right]
  = m_K c\,\frac{1}{(1-\beta^2)^{3/2}}
  = m_K c\,\gamma^3.
  \label{eq:dp_dbeta}
\end{equation}
Dividing by $p = m_K c\,\gamma\beta$ to form the relative error:
\begin{equation}
  \boxed{
  \frac{\delta p}{p} = \gamma^2\,\frac{\delta\beta}{\beta}.
  }
  \label{eq:gamma2_amplification}
\end{equation}
This is the central result.  The Lorentz factor squared $\gamma^2$ multiplies any
fractional error in $\beta$ to produce the fractional error in the reconstructed momentum.

\subsubsection{Physical Origin of the Amplification}

Near $\beta \to 1$ the function $p(\beta)$ becomes extremely steep
(Fig.~\ref{fig:p_vs_beta}).  A small absolute error $|\delta\beta|$ therefore maps to
an arbitrarily large absolute error in $p$.  Equivalently, a fixed relative tracking
uncertainty $\delta\beta/\beta$ is amplified by $\gamma^2$, which itself grows
without bound as $\beta \to 1$.

The amplification arises because the kaon is \emph{ultra-relativistic}: its total
energy $E \approx pc$, so a small fractional change in $\beta$ corresponds to a
small fractional change in $E$ but — through the subtraction
$p = \sqrt{E^2 - m_K^2 c^4}/c$ — to a much larger fractional change in $p$ whenever
$m_K c^2 \ll E$.

\subsubsection{Numerical Example — Event 891}

For the first clean reconstructed event in the
$T_{1}\text{-}400$ configuration simulation:

\begin{center}
\begin{tabular}{lrrl}
  \toprule
  Quantity & Value & Unit & Notes \\
  \midrule
  $p_{\text{true}}$  & 7.455  & GeV/$c$ & Monte~Carlo truth \\
  $\beta_{\text{true}}$ & 0.99773 & --- & from truth track \\
  $\gamma_{\text{true}}$ & 15.01  & --- & \\
  $\gamma^2_{\text{true}}$ & 225 & --- & amplification factor \\[4pt]
  $p_{\text{reco}}$  & 4.468  & GeV/$c$ & PoCA reconstruction \\
  $\delta p/p$       & $-40\%$ & --- & large underestimate \\[4pt]
  Implied $\delta\beta/\beta$ & $-0.18\%$ & --- &
        from Eq.~\eqref{eq:gamma2_amplification}: $-0.40/225$ \\
  \bottomrule
\end{tabular}
\end{center}

A sub-$0.2\%$ error in $\beta$ — well within detector resolution — produces a
$40\%$ underestimate of the momentum because $\gamma^2 = 225$.
Reducing tracking resolution cannot eliminate this effect; it can only marginally
reduce $\delta\beta$ while the $\gamma^2$ factor remains large.

\subsubsection{Identifying the Source of $\delta\beta$}

The measured $\beta$ depends on two independently reconstructed quantities:

\begin{equation}
  \beta_{\text{reco}} = \frac{L_{\text{reco}}}{c\,\Delta t_{\text{reco}}},
  \label{eq:beta_components}
\end{equation}

where
\begin{equation}
  L_{\text{reco}} = \bigl|\bm{x}_{\text{PoCA}} - \bm{x}_{\text{production}}\bigr|
  \label{eq:L_reco}
\end{equation}
is the flight-path length obtained from the PoCA vertex position $\bm{x}_{\text{PoCA}}$,
and $\Delta t_{\text{reco}}$ is derived from pion TOF measurements combined with
a back-propagation along the fitted pion tracks to the kaon production point.

Errors in either component contribute to $\delta\beta$:
\begin{equation}
  \frac{\delta\beta}{\beta} \approx \frac{\delta L}{L} - \frac{\delta(\Delta t)}{\Delta t}.
  \label{eq:beta_error_decomp}
\end{equation}

The PoCA algorithm minimises the separation between two straight-line tracks.
Even when the separation residual is small (e.g.\ $|\delta_{\text{PoCA}}| < 0.1\;\text{cm}$
for Event~891), the reconstructed $z$-coordinate of the vertex is sensitive to the
ratio of track slopes, and any bias in $z$ propagates directly into $L$ and hence
into $\beta$.

\subsubsection{Why Improving Track Resolution Helps Only Marginally}

Let $\sigma_\beta$ be the standard deviation of $\beta_{\text{reco}}$ arising from
finite track resolution.  The standard deviation of the reconstructed momentum is then:
\begin{equation}
  \sigma_p \approx m_K c\,\gamma^3\,\sigma_\beta
             = p\,\gamma^2\,\frac{\sigma_\beta}{\beta}.
  \label{eq:sigma_p}
\end{equation}
Halving $\sigma_\beta$ (e.g.\ by doubling the number of straw layers) halves
$\sigma_p$ — a modest gain — but the $\gamma^2$ prefactor is unchanged.  At
$p \gtrsim 5\;\text{GeV}/c$ where $\gamma^2 \gtrsim 100$, the achievable momentum
resolution is fundamentally limited by the decay-time method and the kaon mass, not
by the tracker hit resolution alone.

To improve sensitivity at high momentum one would need to either:
\begin{enumerate}
  \item measure the flight path $L$ independently (e.g.\ a displaced vertex tag), or
  \item measure the kaon energy directly (calorimetry).
\end{enumerate}
Neither is currently available in this detector geometry.

\subsubsection{Summary}

\begin{itemize}
  \item The relative momentum error scales as $\gamma^2 \times \delta\beta/\beta$
        (Eq.~\eqref{eq:gamma2_amplification}).
  \item At $p \sim 7.5\;\text{GeV}/c$, $\gamma^2 \approx 225$; a $0.18\%$ $\beta$
        error causes a $40\%$ momentum underestimate.
  \item Tracking resolution improvements reduce $\delta\beta$ but cannot remove
        the $\gamma^2$ amplification.
  \item The dominant error source — $z$-vertex bias or timing-chain bias — is
        isolated by the diagnostic reconstruction branches described in
        Section~\ref{sec:diagnostic_branches}.
\end{itemize}

%
%  Required packages in parent preamble:
%    \usepackage{amsmath, amssymb, booktabs, array, bm}
%
%  Describes the fundamental gamma^2 amplification systematic
%  affecting kaon momentum reconstruction from decay-time beta
%  measurement in the KLong simulation.
% ============================================================

\subsection{Systematic Momentum Underestimation: the $\gamma^2$ Amplification Problem}

\subsubsection{The Momentum Formula}

The kaon momentum is reconstructed by measuring its velocity $\beta = v/c$ from the
flight path $L$ and the kaon decay time $\Delta t$:
\begin{equation}
  \beta_{\text{reco}} = \frac{L}{c\,\Delta t},
  \label{eq:beta_reco}
\end{equation}
and then inverting the relativistic energy–momentum relation to obtain
\begin{equation}
  p = m_K\,\gamma\,\beta\,c = \frac{m_K\,\beta\,c}{\sqrt{1-\beta^2}},
  \label{eq:p_from_beta}
\end{equation}
where $m_K = 497.6\;\text{MeV}/c^2$ is the kaon mass.

\subsubsection{Error Propagation and the $\gamma^2$ Factor}

Differentiating Eq.~\eqref{eq:p_from_beta} with respect to $\beta$:
\begin{equation}
  \frac{\mathrm{d}p}{\mathrm{d}\beta}
  = m_K c\,\frac{\mathrm{d}}{\mathrm{d}\beta}\!\left[\frac{\beta}{\sqrt{1-\beta^2}}\right]
  = m_K c\,\frac{1}{(1-\beta^2)^{3/2}}
  = m_K c\,\gamma^3.
  \label{eq:dp_dbeta}
\end{equation}
Dividing by $p = m_K c\,\gamma\beta$ to form the relative error:
\begin{equation}
  \boxed{
  \frac{\delta p}{p} = \gamma^2\,\frac{\delta\beta}{\beta}.
  }
  \label{eq:gamma2_amplification}
\end{equation}
This is the central result.  The Lorentz factor squared $\gamma^2$ multiplies any
fractional error in $\beta$ to produce the fractional error in the reconstructed momentum.

\subsubsection{Physical Origin of the Amplification}

Near $\beta \to 1$ the function $p(\beta)$ becomes extremely steep
(Fig.~\ref{fig:p_vs_beta}).  A small absolute error $|\delta\beta|$ therefore maps to
an arbitrarily large absolute error in $p$.  Equivalently, a fixed relative tracking
uncertainty $\delta\beta/\beta$ is amplified by $\gamma^2$, which itself grows
without bound as $\beta \to 1$.

The amplification arises because the kaon is \emph{ultra-relativistic}: its total
energy $E \approx pc$, so a small fractional change in $\beta$ corresponds to a
small fractional change in $E$ but — through the subtraction
$p = \sqrt{E^2 - m_K^2 c^4}/c$ — to a much larger fractional change in $p$ whenever
$m_K c^2 \ll E$.

\subsubsection{Numerical Example — Event 891}

For the first clean reconstructed event in the
$T_{1}\text{-}400$ configuration simulation:

\begin{center}
\begin{tabular}{lrrl}
  \toprule
  Quantity & Value & Unit & Notes \\
  \midrule
  $p_{\text{true}}$  & 7.455  & GeV/$c$ & Monte~Carlo truth \\
  $\beta_{\text{true}}$ & 0.99773 & --- & from truth track \\
  $\gamma_{\text{true}}$ & 15.01  & --- & \\
  $\gamma^2_{\text{true}}$ & 225 & --- & amplification factor \\[4pt]
  $p_{\text{reco}}$  & 4.468  & GeV/$c$ & PoCA reconstruction \\
  $\delta p/p$       & $-40\%$ & --- & large underestimate \\[4pt]
  Implied $\delta\beta/\beta$ & $-0.18\%$ & --- &
        from Eq.~\eqref{eq:gamma2_amplification}: $-0.40/225$ \\
  \bottomrule
\end{tabular}
\end{center}

A sub-$0.2\%$ error in $\beta$ — well within detector resolution — produces a
$40\%$ underestimate of the momentum because $\gamma^2 = 225$.
Reducing tracking resolution cannot eliminate this effect; it can only marginally
reduce $\delta\beta$ while the $\gamma^2$ factor remains large.

\subsubsection{Identifying the Source of $\delta\beta$}

The measured $\beta$ depends on two independently reconstructed quantities:

\begin{equation}
  \beta_{\text{reco}} = \frac{L_{\text{reco}}}{c\,\Delta t_{\text{reco}}},
  \label{eq:beta_components}
\end{equation}

where
\begin{equation}
  L_{\text{reco}} = \bigl|\bm{x}_{\text{PoCA}} - \bm{x}_{\text{production}}\bigr|
  \label{eq:L_reco}
\end{equation}
is the flight-path length obtained from the PoCA vertex position $\bm{x}_{\text{PoCA}}$,
and $\Delta t_{\text{reco}}$ is derived from pion TOF measurements combined with
a back-propagation along the fitted pion tracks to the kaon production point.

Errors in either component contribute to $\delta\beta$:
\begin{equation}
  \frac{\delta\beta}{\beta} \approx \frac{\delta L}{L} - \frac{\delta(\Delta t)}{\Delta t}.
  \label{eq:beta_error_decomp}
\end{equation}

The PoCA algorithm minimises the separation between two straight-line tracks.
Even when the separation residual is small (e.g.\ $|\delta_{\text{PoCA}}| < 0.1\;\text{cm}$
for Event~891), the reconstructed $z$-coordinate of the vertex is sensitive to the
ratio of track slopes, and any bias in $z$ propagates directly into $L$ and hence
into $\beta$.

\subsubsection{Why Improving Track Resolution Helps Only Marginally}

Let $\sigma_\beta$ be the standard deviation of $\beta_{\text{reco}}$ arising from
finite track resolution.  The standard deviation of the reconstructed momentum is then:
\begin{equation}
  \sigma_p \approx m_K c\,\gamma^3\,\sigma_\beta
             = p\,\gamma^2\,\frac{\sigma_\beta}{\beta}.
  \label{eq:sigma_p}
\end{equation}
Halving $\sigma_\beta$ (e.g.\ by doubling the number of straw layers) halves
$\sigma_p$ — a modest gain — but the $\gamma^2$ prefactor is unchanged.  At
$p \gtrsim 5\;\text{GeV}/c$ where $\gamma^2 \gtrsim 100$, the achievable momentum
resolution is fundamentally limited by the decay-time method and the kaon mass, not
by the tracker hit resolution alone.

To improve sensitivity at high momentum one would need to either:
\begin{enumerate}
  \item measure the flight path $L$ independently (e.g.\ a displaced vertex tag), or
  \item measure the kaon energy directly (calorimetry).
\end{enumerate}
Neither is currently available in this detector geometry.

\subsubsection{Summary}

\begin{itemize}
  \item The relative momentum error scales as $\gamma^2 \times \delta\beta/\beta$
        (Eq.~\eqref{eq:gamma2_amplification}).
  \item At $p \sim 7.5\;\text{GeV}/c$, $\gamma^2 \approx 225$; a $0.18\%$ $\beta$
        error causes a $40\%$ momentum underestimate.
  \item Tracking resolution improvements reduce $\delta\beta$ but cannot remove
        the $\gamma^2$ amplification.
  \item The dominant error source — $z$-vertex bias or timing-chain bias — is
        isolated by the diagnostic reconstruction branches described in
        Section~\ref{sec:diagnostic_branches}.
\end{itemize}
